\abstract{
\thispagestyle{fancy}

%\begin{otherlanguage}{ukrainian}
%2. Для ознайомлення зі змістом та результатами дисертації подається державною та англійською мовами анотація - узагальнений короткий виклад її основного змісту. В анотації дисертації мають бути стисло представлені основні результати дослідження із зазначенням наукової новизни та за наявності практичного значення.

%В анотації також вказуються:

%прізвище та ініціали здобувача;

%назва дисертації;

%вид дисертації та науковий ступінь, на який претендує здобувач;

%спеціальність (шифр і назва);

%найменування вищого навчального закладу або найменування наукової установи, у якому (якій) здійснювалася підготовка;

%найменування наукової установи або найменування вищого навчального закладу, у спеціалізованій вченій раді якої (якого) відбудеться захист;

%місто, рік.

%Зразок анотації наведено у додатку 2 до цих Вимог

%Обсяг анотації становить 0,2 - 0,3 авторських аркуша.

%Анотація може подаватися також третьою мовою, пов’язаною з предметом дослідження.

%3. Наприкінці анотації наводяться ключові слова відповідною мовою. Сукупність ключових слів повинна відповідати основному змісту наукової праці, відображати тематику дослідження і забезпечувати тематичний пошук роботи. Кількість ключових слів становить від п’яти до п’ятнадцяти. Ключові слова подають у називному відмінку, друкують в рядок через кому.

%4. Після ключових слів наводиться список публікацій здобувача за темою дисертації. Вказуються наукові праці:

%в яких опубліковані основні наукові результати дисертації;

%які засвідчують апробацію матеріалів дисертації;

%які додатково відображають наукові результати дисертації.

%\end{otherlanguage}

\begin{otherlanguage}{ukrainian}

\textit{Фадєєва Л.О.} Методика навчання цифрової історії майбутніх фахівців. -- Кваліфікаційна наукова праця на правах рукопису.

Дисертація на здобуття наукового ступеня доктора філософії за спеціальністю 011 Освітні, педагогічні науки. -- Криворізький державний педагогічний університет, Кривий Ріг, 2024.

{У ході розв'язання наукової проблеми підготовки майбутніх викладачів \linebreak STEM-дисциплін до застосування технологій доповненої реальності були  отримані наступні результати.

Бібліометричний аналіз джерел в галузі STEM та підготовки викладачів ідентифікував 21 ключове поняття, систематизовані у чотири кластери: ``STEM-освіта і підготовка кадрів'', ``Професійна підготовка вчителів і початкова освіта'', ``Опитування щодо STEM-освіти'' та ``Електронне навчання і обчислювальне мислення в підготовці майбутніх викладачів STEM-дисциплін''.

За результатами систематичного аналізу було встановлено, що: а) 
    STEM-освіта позитивно впливає на досягнення у природничих науках і математиці, але існує дефіцит фахівців у цих галузях;
б)    зниження інтересу до STEM-дисциплін~-- глобальна проблема, яку можна вирішити за допомогою привабливих та доступних освітніх програм;
    в) нерівність у сфері STEM -- глобальна проблема, яку можна вирішити через створення STEM-лабораторій та підготовку вчителів з урахуванням гендерних аспектів;
    г) обчислювальне мислення -- важлива складова STEM-освіти, що може бути впроваджена через семінари для вчителів, онлайн-курси та методичну підтримку;
    д) професійна підготовка та соціальний статус учителів є стратегічно важливими для STEM-освіти;
    е) зацікавлення у STEM-дис\-цип\-лі\-нах може бути підвищено за допомогою ІКТ, зокрема доповненої реальності, віртуальної реальності та робототехніки.

Отримані результати дають змогу запропонувати такі рекомендації для поліпшення STEM-освіти:
а) інтегрувати STEM-підхід у програми підготовки вчителів;
б) розвивати проєктне мислення, цифрові та STEM-навички у вчителів та учнів;
в) запроваджувати міждисциплінарні STEM-проєкти;
г) використовувати активні та практико орієнтовані методи навчання;
д) збільшувати доступність STEM-освіти для всіх учнів;
е) створювати STEM-лабораторії у закладах загальної середньої освіти;
ж) підвищувати соціальний статус викладачів STEM-дисциплін;
з) запроваджувати засоби ІКТ у навчання STEM-дисциплін.

Педагогічні умови включають матеріальні, методичні, організаційні та інші фактори, що забезпечують ефективність підготовки майбутніх викладачів STEM-дисциплін до використання доповненої реальності у своїй професійній діяльності. Опитування 94 респондентів, переважно викладачів STEM-дисциплін, що використовують доповнену реальність, виявило наступні педагогічні умови:
1)
    забезпечення доступності мобільних апаратних засобів доповненої реальності та імерсивних цифрових освітніх ресурсів майбутнім викладачам STEM-дисциплін;
 2)   уведення до змісту підготовки питань, пов'язаних із використанням доповненої реальності у навчанні STEM-дисциплiн;
    3) застосування дослідницького підходу та інтерактивних технологій у процесі підготовки майбутніх викладачiв STEM-дисциплiн;
    4) набуття практичного досвіду застосування технологій доповненої реальності у навчанні STEM-дисциплiн.

Було встановлено, що доступність мобільних пристроїв для доповненої реальності та імерсивних цифрових освітніх ресурсів для майбутніх викладачів \linebreak STEM-дисциплін забезпечується частково. Через обмежену кількість імерсивних цифрових освітніх ресурсів майбутні викладачі STEM-дисциплін повинні брати участь у їх розробці, що сприяє їх професійному розвитку.

Для впровадження питань, пов'язаних із застосуванням доповненої реальності у навчанні STEM-дисциплін, було розроблено елементи методики навчання майбутніх викладачів STEM-дисциплін створення імерсивних освітніх ресурсів.

Застосування дослідницького підходу та інтерактивних технологій у процесі підготовки майбутніх викладачів STEM-дисциплін передбачало два напрями: створення STEM-проєктів із доповненою реальністю та створення системи завдань, що сприяють пошуковій та творчій активності студентів. Такі підходи допомагають підвищити інтерактивність та ефективність навчання STEM-дисциплін.

У процесі дослідно-експериментальної роботи було створено три версії навчального курсу для майбутніх викладачів STEM-дисциплін, що сприяє позитивній динаміці в підготовці майбутніх викладачів до використання імерсивних освітніх ресурсів. Більшість учасників експерименту планують використовувати доповнену реальність у своїй професійній діяльності, що свідчить про позитивну ефективність навчання.



Наукова новизна отриманих результатів полягає в тому, що
    {вперше} 
    виокремлено та теоретично обґрунтовано
    педагогічні умови підготовки майбутніх викладачів STEM-дисцип\-лін до застосування технологій доповненої реальності у професійній діяльності: а)
   забезпечення доступності мобільних апаратних засобів доповненої реальності та імерсивних цифрових освітніх ресурсів майбутнім викладачам STEM-дисциплін; 
        б)
   уведення до змісту підготовки питань, пов'язаних iз використанням доповненої реальності у навчанні STEM-дисциплін; 
        в)
   застосування дослідницького підходу та інтерактивних технологій у процесі підготовки майбутніх викладачів STEM-дис\-цип\-лiн; 
        г)
   набуття практичного досвіду застосування технологій доповненої реальності у навчанні STEM-дисциплін;
    {удосконалено} зміст професійної підготовки майбутніх викладачів STEM-дис\-цип\-лін;
    {набули подальшого розвитку} теорія та методика професійної підготовки майбутніх викладачів STEM-дисциплін.
    

{Практичне значення отриманих результатів} полягає в тому, що розроблено окремі елементи методики навчання майбутніх викладачів STEM-дис\-цип\-лін створення імерсивних освітніх ресурсів у складі  електронного навчального курсу та посібника до нього. 



\textit{Ключові слова}: STEM-освіта, STEM-дисципліни, професійна підготовка викладачів, педагогічні умови,  методика навчання, доповнена реальність, віртуальна реальність, імерсивні технології, цифрові освітні ресурси, імерсивні освітні ресурси, мобільні технології.}

\end{otherlanguage}

% \fi
% \newpage
% \iffalse

\begin{center}
\large
    \textbf{ABSTRACT}
\end{center}

\textit{Fadieieva L. O.} Methodology of using information and cognitive technologies of adaptive learning of pedagogical universities students. -- Qualifying scientific work on manuscript rights.

Dissertation for the Doctor of Philosophy degree in the speciality 011 Education and Pedagogical Sciences. -- Kryvyi Rih State Pedagogical University, Kryvyi Rih, 2024.

{In the course of solving the scientific problem of preparing future STEM teachers to use augmented reality technologies, the following results were obtained. 

Bibliometric analysis of STEM and teacher training sources identified 21 key concepts, systematised into four clusters: “STEM education and personnel training”, “Tea\-cher training and primary education”, “STEM education surveys”, and “E-learning and computational thinking in training future STEM teachers”.

The systematic analysis found that: a) 
STEM education positively impacts achievement in science and maths, but there is a shortage of specialists in these fields;
b) the decline of interest in STEM subjects is a global problem that can be solved through attractive and accessible educational programs;
c) inequality in STEM is a global problem that can be solved by creating STEM labs and training teachers with gender considerations;
d) computational thinking is an essential component of STEM education that can be implemented through teacher workshops, online courses and methodological support;
e) professional training and social status of teachers are strategically crucial for STEM education;
f) interest in STEM subjects can be increased through ICT, specifically augmented reality, virtual reality and robotics.  

The obtained results make it possible to propose the following recommendations for improving STEM education:
a) integrate the STEM approach into teacher training programmes; 
b) develop project thinking, digital and STEM skills in teachers and pupils;
c) implement interdisciplinary STEM projects;
d) use active and practice-oriented teaching methods;
e) increase the accessibility of STEM education for all pupils;
f) create STEM labs in secondary schools;
g) raise the social status of STEM teachers;
h) implement ICT tools in STEM teaching.

Pedagogical conditions include material, methodological, organisational and other factors that ensure the effectiveness of preparing future STEM teachers to use AR in their professional activities. A survey of 94 respondents, mostly STEM teachers using augmented reality, identified the following pedagogical conditions:
1) providing accessibility of mobile augmented reality hardware and immersive digital educational resources to future STEM teachers; 
2) introducing content on the use of augmented reality in STEM teaching into the training;
3) applying a research approach and interactive technologies in the training of future STEM teachers; 
4) gaining practical experience in using augmented reality technologies in STEM teaching.

It was found that the accessibility of mobile AR devices and immersive digital educational resources for future STEM teachers is only partially ensured. Due to the limited number of immersive digital educational resources, future STEM teachers should participate in their development, contributing to their professional growth.

To introduce content on using augmented reality in STEM teaching, elements of a methodology for teaching future STEM teachers to create immersive educational resources were developed. 

Applying a research approach and interactive technologies in the training of future STEM teachers involved two areas: creating AR-enhanced STEM projects and developing a system of tasks that promote inquiry-based and creative activity of students. Such approaches help increase the interactivity and effectiveness of STEM teaching.

In the process of research and experimental work, three versions of a training course for future STEM teachers were created, contributing to positive dynamics in preparing future teachers to use immersive educational resources. Most participants in the experiment plan to use augmented reality in their professional activities, indicating the positive effectiveness of the training.

The scientific novelty of the obtained results lies in the fact that for the first time:} 
\begin{itemize}
    \item 
{the pedagogical conditions for preparing future STEM teachers to use augmented reality technologies in professional activities were identified and theoretically substantiated: a) providing accessibility of mobile AR hardware and immersive digital educational resources to future STEM teachers; b) introducing content on the use of augmented reality in STEM teaching into the training; c) applying a research approach and interactive technologies in the training of future STEM teachers; d) gaining practical experience in using augmented reality technologies in STEM teaching;}
    \item 
{the content of professional training for future STEM teachers has been improved; }
    \item 
{the theory and professional training methods for future STEM teachers have been further developed.}
\end{itemize}

{The practical significance of the obtained results lies in that individual elements of the methodology for teaching future STEM teachers to create immersive educational resources have been developed as part of an e-learning course and accompanying textbook.

\textit{Keywords}: STEM education, STEM disciplines, teachers' professional training, pedagogical conditions, teaching methods, augmented reality, virtual reality, immersive technologies, digital educational resources, immersive educational resources, mobile technologies.}


\begin{center}
    \textbf{List of education seeker publications}
\end{center}


\printpratsi

} % Write your abstract between the {}
